\startchapter{Related Literature}
\label{chapter:related_literature}

This chapter reviews the literature underpinning the methods evaluated in this thesis, situating them within their respective disciplinary traditions before drawing them together around the central problem of PHC plume reconstruction. Section \ref{sec:spat_int_earthSciences} traces the historical development of spatial interpolation in the earth sciences, from the practical origins of IDW to the more statistically grounded Kriging framework, establishing the conceptual foundation for the distance-based methods benchmarked in this work. Section \ref{sec:comp_int_methods} surveys comparative studies of these interpolation methods across meteorology, soil science, and contaminant hydrology, highlighting a persistent gap. No existing study evaluates spatial interpolation alongside parameter estimation techniques within a unified comparative framework. Section \ref{sec:param_est_soil_groundwater} turns to parameter estimation as it has developed in groundwater and soil modelling, where methods such as PEST have become standard tools for model calibration but have not previously been assessed for their capacity to reconstruct spatial distributions from sparse observations. Finally, Section \ref{sec:fdvar_da} introduces 4D-Var, a variational data assimilation technique developed in numerical weather prediction, and discusses its methodological kinship to parameter estimation despite arising from a distinct disciplinary context. Together, these four bodies of literature reveal that while spatial interpolation and parameter estimation have each been extensively studied and benchmarked within their own domains, they have not been directly compared for the specific task of reconstructing contaminant plumes from sparse sensor networks, which is the gap this thesis aims to fill.

\section{Spatial Interpolation in the Earth Sciences}
\label{sec:spat_int_earthSciences}

The application of spatial interpolation to subsurface and soil characterization has a long history rooted in the practical constraints of field measurement. Early mining surveyors faced significant technological and logistical limitations that precluded dense sampling in either time or space, necessitating intelligent methods to extract maximum information from sparse observational datasets \cite{krigHistory, spatInt_hist}. Inverse Distance Weighting (IDW) had long served as a versatile interpolation tool across numerous disciplines, relying on the intuitive principle that nearby observations carry greater predictive weight than distant ones. However, the specific demands of geostatistical analysis in mining and geology motivated the development of Kriging, a family of methods that more rigorously exploits the spatial correlation structure of the data through a fitted variogram model \cite{krigHistory}. While IDW has remained relatively unchanged in its formulation, the Kriging framework has since expanded into a suite of variants including Ordinary Kriging, Universal Kriging, Co-Kriging, and Indicator Kriging, among others; each tailored to address specific assumptions about data stationarity, secondary covariates, or non-Gaussian behaviour \cite{kriging_overview}. Despite their differing mathematical underpinnings, these methods share the common feature of producing distance-weighted linear combinations of the available measurements.

Several additional spatial interpolation approaches exist beyond the distance-weighting paradigm. These include spline-based methods, radial basis functions, and global or local polynomial regression techniques \cite{mitas1999spatial}. Although these methods appear in the broader interpolation literature and are evaluated in several of the comparative studies discussed below, they were not selected for evaluation to simplify the comparison and incorporate the addition of parameter estimation methods.

\section{Comparative Studies of Interpolation Methods}
\label{sec:comp_int_methods}

A considerable body of literature exists evaluating and comparing spatial interpolation methods across a range of earth science disciplines. In meteorology and climatology, IDW, Kriging, and related approaches have been benchmarked extensively for the interpolation of precipitation fields and other weather variables across spatial domains ranging from individual cities to continents \cite{rainWater_spatComp, rainWater_spatComp2}. These studies generally reveal that the relative performance of methods is context-dependent, varying with the spatial autocorrelation structure of the variable of interest, the density and spatial distribution of the monitoring network, and the terrain complexity of the study domain.

In soil science, these interpolation methods have been applied and compared on a variety of physical and chemical soil properties, including moisture content, temperature, and nutrient concentrations \cite{spatInt_drought, spatInt_soilProp}. Kriging-based methods have frequently demonstrated advantages over simpler approaches in these settings, particularly where underlying spatial structure can be characterized through variogram analysis. More directly relevant to the present work, several studies have extended these comparisons to contaminant distributions in soil and shallow subsurface environments, including heavy metal enrichment in mine-affected soils and petroleum hydrocarbon (PHC) contamination \cite{spatInt_mines, spatInt_PHC}. Collectively, this body of work establishes a solid baseline for the performance of classical interpolation methods in soil contamination contexts. None of these studies evaluate spatial interpolation methods alongside parameter estimation techniques in a unified framework which is a gap that this thesis attempts to address. 

\section{Parameter Estimation in Soil and Groundwater Modelling}
\label{sec:param_est_soil_groundwater}

The application of formal parameter estimation techniques to subsurface flow and transport modelling has developed largely in parallel with the spatial interpolation literature. Early efforts focused on groundwater flow modelling and the characterization of flow dynamics within the vadose zone, where the governing equations are well-established but parameterized by spatially variable quantities that are difficult to measure directly \cite{parEst_first}. These early methods have since been codified and made broadly accessible to the soil and groundwater modelling community through standardized software libraries, most notably the PEST (Parameter ESTimation) suite and its extensions \cite{pestManual}. Today, parameter estimation tools are employed in the calibration of regional groundwater models for municipal water management and the simulation of water movement through heterogeneous soil profiles \cite{parEst_ex1, parEst_ex2}.

Comparative assessments of parameter estimation approaches have been conducted primarily in the context of groundwater flow, with particular attention to the estimation of hydraulic conductivity fields and the calibration of regional flow models \cite{parEst_comp_groundwaterFlow, parEst_comp_hydraulicConductivity}. The transport of contaminants through soil is a closely related but comparatively less-studied application, in part because contaminant fate-and-transport models typically build upon and inherit the structural assumptions of the underlying groundwater flow models. While parameter estimation methods have been thoroughly evaluated for their utility in model calibration, they have not previously been assessed within a spatial interpolation framework, evaluating their capacity to reproduce observed spatial distributions of a variable of interest across a monitoring network. This thesis presents, to the author's knowledge, the first systematic comparison of parameter estimation methods against classical spatial interpolation techniques for this purpose.

\section{Four-Dimensional Variational Data Assimilation}
\label{sec:fdvar_da}

Four-Dimensional Variational Data Assimilation (4D-Var) originates in the field of numerical weather prediction (NWP), where it was developed as a rigorous mathematical framework for optimally combining sparse atmospheric observations with a high-fidelity dynamical model of the atmosphere \cite{da_history}. By minimizing a cost function that penalizes both the misfit between model predictions and observations and deviations from a background state, 4D-Var produces the best estimate of the atmospheric state. This state typically encompasses variables such as temperature, pressure, humidity, and wind at a given time, consistent with both the observations and the model dynamics \cite{da_history, da_overview}. The method is mathematically related to the parameter estimation approaches described in the preceding section, sharing the same variational optimization framework, though it differs in its application to time-evolving dynamical systems and in the particular structure of its adjoint-based gradient computation.

Given the computational cost associated with running and differentiating through state-of-the-art NWP models, operational 4D-Var implementations have historically been constrained to a limited number of model evaluations. The method has demonstrated value beyond the NWP context, including applications to the spatial interpolation problem \cite{fdvar_interpolation}. The decreasing cost of high-performance computing has further stimulated comparative research, and a growing number of studies have benchmarked 4D-Var against ensemble-based data assimilation methods \cite{fdvar_vs_ensemble,lorenc2003potential,bannister2017review}. While 4D-Var has not yet been formally adopted within the soil parameter estimation community, its conceptual proximity to parameter estimation and its demonstrated capacity for spatial state reconstruction make it a relevant reference point for the methodological comparisons undertaken in this thesis.

\section{Summary}
\label{sec:lit_review_summary}

Taken together, this body of literature establishes a clear picture of two methodological traditions that have developed largely in isolation from one another. Spatial interpolation methods, from IDW through the various Kriging formulations, have been extensively benchmarked within the earth sciences, including on soil contaminant distributions closely related to the PHC plumes considered in this thesis. Parameter estimation methods, by contrast, have matured within the groundwater and soil modelling community as tools for model calibration, while 4D-Var has been developed and refined within numerical weather prediction as a means of reconstructing spatial states from sparse observations. Despite their shared reliance on optimization against observed data, and despite each tradition individually demonstrating value for problems structurally similar to PHC plume reconstruction, no existing study has directly compared spatial interpolation against parameter estimation for this task. Given the practical importance of accurately reconstructing contaminant plumes from the sparse sensor networks typical of active remediation sites, and given that parameter estimation methods have not previously been assessed for this specific capacity, a systematic comparison of these two methodological traditions is necessary to identify the most effective approach for PHC plume reconstruction. This thesis undertakes that comparison.